\section{Actual formal division by internally constructed vertex series}

Keep the ACTUAL whole formal ideal $I\subset S=\mathbb C[[X_1,\ldots,X_n]]$,
its genuine initial diagram $\mathcal N$, its internally constructed finite
vertex set $V$, and the internally chosen genuine monic series $f_a\in I$
with lowest exponent $a$, for $a\in V$. All products are the genuine
finite-antidiagonal formal products. The construction below is a direct
coefficient version of formal division in
\cite[Theorem~3.17, Lemma~7.3, pp.~261--264]{bierstone-milman97-invent-math}.
We allow any fixed internally chosen owner of each exponent in the diagram;
this refines the same vertex cones into disjoint parts. The usual first-vertex
partition is one such choice. This harmless partition convention does not
assume or modify the whole ideal or the actual series ring.

\begin{lemma}[Actual vertex owners and well founded genuine coefficient recursion]
\label{hir:actual-formal-vertex-owner-coefficient-recursion}
\uses{hir:actual-formal-initial-diagram-finite-monic-vertices,hir:actual-formal-initial-exponent-coefficient-product}
Internally choose for each $c\in\mathcal N$ one actual $o(c)\in V$ with
$o(c)\le c$; for $c\notin\mathcal N$, put $o(c)=\mathrm{none}$.
For every actual $F\in S$ there is a genuinely well defined complex
coefficient function $t_F(c)$ given by
\[
 t_F(c)=[X^c]F-
 \sum_{a\in V}\sum_{\beta+\gamma=c\atop\beta+a\prec c,\ o(\beta+a)=a}
 t_F(\beta+a)[X^\gamma]f_a .
\]
Every displayed inner sum is the actual finite exponent antidiagonal;
every recursive call is strictly smaller in the genuine well monomial
order. Neither the owner nor a coefficient solution is a caller premise.
\lean{Hironaka.HironakaActualFormalVertexDivision.exists_actualDivisionOwner,Hironaka.HironakaActualFormalVertexDivision.actualDivisionOwner,Hironaka.HironakaActualFormalVertexDivision.actualDivisionOwner_le,Hironaka.HironakaActualFormalVertexDivision.actualDivisionOwner_eq_none_iff,Hironaka.HironakaActualFormalVertexDivision.actualVertexSeries,Hironaka.HironakaActualFormalVertexDivision.actualVertexSeries_spec,Hironaka.HironakaActualFormalVertexDivision.actualVertexSeries_coefficient,Hironaka.HironakaActualFormalVertexDivision.actualVertexSeries_coefficient_lower,Hironaka.HironakaActualFormalVertexDivision.actualDivisionCoefficient,Hironaka.HironakaActualFormalVertexDivision.actualDivisionTail,Hironaka.HironakaActualFormalVertexDivision.actualDivisionCoefficient_eq}
\leanok
\end{lemma}
\begin{proof}
The actual finite-vertex coverage theorem gives an actual dominating
vertex for every diagram exponent, so choice constructs the stated
owner, with the stated pointwise inequality. Outside the actual diagram
there is no owner. Apply well founded recursion to the actual well
monomial exponent order. At stage $c$, the formula uses only the finitely
many actual antidiagonal pairs and only the already defined values at
$\beta+a\prec c$. Thus it gives an actual complex number at every
exponent, with the displayed defining equation. No infinite numerical
sum or convergence assertion enters the construction.
\end{proof}

\begin{theorem}[Actual formal quotients and remainder have the required disjoint supports]
\label{hir:actual-formal-vertex-division-supports}
\uses{hir:actual-formal-vertex-owner-coefficient-recursion}
Define GENUINE actual formal series $Q_a(F),R(F)\in S$ by their coefficients:
\[
 [X^\beta]Q_a(F)=
 \begin{cases}t_F(\beta+a)&o(\beta+a)=a,\\0&\text{otherwise},\end{cases}
 \qquad
 [X^c]R(F)=
 \begin{cases}t_F(c)&c\notin\mathcal N,\\0&c\in\mathcal N.\end{cases}
\]
The shifted supports $a+\operatorname{supp}Q_a(F)$ lie in the actual
disjoint owner parts of $\mathcal N$, and
$\operatorname{supp}R(F)\cap\mathcal N=\varnothing$.
\lean{Hironaka.HironakaActualFormalVertexDivision.actualFormalQuotient,Hironaka.HironakaActualFormalVertexDivision.actualFormalRemainder,Hironaka.HironakaActualFormalVertexDivision.actualFormalQuotient_support,Hironaka.HironakaActualFormalVertexDivision.actualFormalRemainder_coefficient_diagram}
\leanok
\end{theorem}
\begin{proof}
An actual formal series is an arbitrary complex coefficient function on
the actual exponent set, with its existing genuine ring operations.
The displayed functions therefore define actual series directly.
A nonzero quotient coefficient forces its defining owner condition;
that condition places the shifted exponent in the true diagram. Distinct
owners cannot have the same shifted exponent because an actual owner
is a single value. A remainder coefficient at a diagram exponent is
literally zero by its definition, proving the last assertion.
\end{proof}

\begin{theorem}[The actual finite sum of genuine formal products equals the original series]
\label{hir:actual-formal-vertex-division-identity}
\uses{hir:actual-formal-vertex-division-supports,hir:actual-formal-initial-exponent-coefficient-product}
For EVERY actual $F\in S$, the internally constructed series satisfy the
WHOLE genuine formal equality
\[
 F=\sum_{a\in V} Q_a(F)f_a+R(F).
\]
\lean{Hironaka.HironakaActualFormalVertexDivision.actual_formal_vertex_division}
\leanok
\end{theorem}
\begin{proof}
Fix an actual exponent $c$ and expand the coefficient of the displayed
finite sum using the genuine product antidiagonals. The terms with
$\beta+a\prec c$ are exactly the correction sum in the defining
recursion, by the actual quotient coefficient formula.
For every remaining antidiagonal pair, a nonzero coefficient of $f_a$
forces $a\preceq\gamma$. If $\beta+a\not\prec c=\beta+\gamma$,
ordered cancellation then forces $\gamma=a$. Such a term exists only
when $a\le c$ pointwise, has $\beta=c-a$, and contributes $t_F(c)$
exactly when $o(c)=a$, since $f_a$ has actual coefficient $1$ there.
Thus if $c\in\mathcal N$, exactly its one genuine owner contributes
that remaining $t_F(c)$ and the remainder coefficient is zero. If
$c\notin\mathcal N$, no owner contributes and the actual remainder
coefficient is $t_F(c)$. In both cases the coefficient of the whole right
side is the correction plus $t_F(c)$, which is precisely $[X^c]F$
by the true recursion equation. Equality of EVERY actual coefficient
proves the actual formal ring equality.
\end{proof}

\begin{theorem}[Internally constructed actual vertex series generate the whole original formal ideal]
\label{hir:actual-formal-vertex-series-whole-ideal-generation}
\uses{hir:actual-formal-vertex-division-identity,hir:actual-formal-ideal-initial-diagram-upper}
For every actual $F\in I$, its internally constructed remainder is zero.
Consequently $I$ is generated by the internally constructed finite actual
family $(f_a)_{a\in V}$, with the genuine whole-series coefficient
decomposition of the preceding theorem. This includes the zero and unit
ideals and requires no basis, quotient, or division certificate.
\lean{Hironaka.HironakaActualFormalVertexDivision.actualFormalRemainder_eq_zero_of_mem,Hironaka.HironakaActualFormalVertexDivision.actual_formal_ideal_member_division,Hironaka.HironakaActualFormalVertexDivision.actual_formal_ideal_eq_span_vertex_series}
\leanok
\end{theorem}
\begin{proof}
Each actual $f_a$ belongs to the SAME whole ideal $I$, so the actual finite
sum $\sum Q_a(F)f_a$ belongs to $I$. The genuine identity therefore
puts $R(F)$ in $I$. If this actual remainder were nonzero, its actual
lowest nonzero coefficient exponent would belong to $\mathcal N$
by the definition of that diagram. But the actual remainder has zero
coefficient at every diagram exponent, contradicting its genuine
lowest nonzero coefficient. Thus $R(F)=0$. The resulting identity puts
every actual member of $I$ in the span of the actual $f_a$; the converse
inclusion follows from their actual membership in $I$. This proves BOTH
whole ideal inclusions internally.
\end{proof}
